\chapter{Le marché des changes}\label{ch:m2:the-fx-market}

En avril 2025, les banques et les dealers du monde entier ont échangé, en moyenne par jour, USD~9.6
mille milliards de devises, soit 28~\% de plus que trois ans plus tôt. Il n'existe
aucune bourse où cela se produit, aucun prix unique et aucune heure de clôture. Le
marché est un réseau de dealers qui cotent entre eux et auprès de leurs clients, quelques
carnets d'ordres électroniques que la plupart des acteurs du monde observent, des dizaines de plateformes
qui rivalisent pour les relier, et un système de règlement-livraison qui veille à ce que chaque partie
d'une transaction livre sa devise. Le dollar américain figure d'un côté de presque neuf
transactions sur dix. Ce chapitre décrit les conventions du marché, sa
structure, les plateformes où il se négocie et ce que les banques centrales en disent dans leur enquête~;
les chapitres suivants en passent les instruments en revue un à un.

\section{Devises, paires et conventions}

\begin{definition}[Paire de devises, devise de base et devise de cotation, pip]\label{def:m2:the-fx-market:pair}
Une \emph{paire de devises}\index{paire de devises} s'écrit au moyen de deux codes de trois lettres,
BASEQUOTE, et son cours $S$ est le nombre d'unités de la seconde, la
\emph{devise de cotation}\index{devise de cotation}, pour une unité de la première, la
\emph{devise de base}\index{devise de base}. Un \emph{pip}\index{pip} est
l'unité conventionnelle de variation d'un cours~: 0.0001 pour la plupart des paires, 0.01 pour les paires
cotées en yens.
\end{definition}

EURUSD à 1.1464 signifie 1.1464 dollars par euro~; USDJPY à 156.87 signifie 156.87
yens par dollar. Le choix de la \omterm{def:m2:the-fx-market:pair}{devise de base} relève d'une convention, non d'un choix~:
les taux journaliers de la Réserve fédérale sont exprimés en unités de devise par dollar, sauf
pour l'euro, la livre sterling et les dollars australien et néo-zélandais, qui sont
exprimés en dollars par unité~; la Banque centrale européenne cote chaque taux par euro,
de sorte que EURGBP s'exprime en livres par euro. Un trader qui \emph{achète} EURUSD achète des euros et
vend des dollars~; une cotation de 1.1462/1.1464 signifie que le dealer achète des euros
au cours acheteur, 1.1462, et les vend au cours vendeur, 1.1464, soit une fourchette de deux \omterm{def:m2:the-fx-market:pair}{pips}.

\begin{definition}[Date de valeur spot, cours croisé]\label{def:m2:the-fx-market:spot}
La \emph{date de valeur spot}\index{date de valeur spot} d'une transaction est le jour où
les deux devises sont échangées~: par convention, le deuxième jour ouvré suivant
la transaction, en ne comptant que les jours ouvrés dans les deux devises, et
le premier pour le couple dollar--dollar canadien. Un \emph{cours croisé}\index{cours croisé} est
le taux entre deux devises dont aucune n'est le dollar, comme EURJPY,
négocié directement ou par l'intermédiaire des deux paires en dollars.
\end{definition}

La règle comporte un détail que tout système doit intégrer. Des économistes de la Réserve fédérale
l'ont décrite en 2003 à partir du couple dollar--yen~: une transaction conclue un mardi se règle
le jeudi~; si le mercredi est férié au Japon, le règlement a lieu le vendredi~; mais si le mercredi est
férié aux États-Unis et non au Japon, il reste fixé au jeudi, car un
jour férié américain intermédiaire ne retarde pas une paire en dollars
(\cref{fig:m2:the-fx-market:spotdate}). Une journée prend fin à 17~h à New York~: une
transaction conclue juste après cette heure est datée du lendemain et se règle un jour plus tard. Tout
ce qui va au-delà du spot est un terme, objet du \cref{ch:m2:fx-swaps-forwards-and-the-cross-currency-basis}.

\begin{omfigure}
\centering
\begin{tikzpicture}[font=\small,x=1.9cm,y=0.75cm]
\foreach \d/\n in {0/mar.,1/mer.,2/jeu.,3/ven.} {\node[font=\footnotesize\bfseries] at (\d,0.9) {\n};}
\node[anchor=east,font=\footnotesize] at (-0.45,0) {aucun jour férié};
\node[anchor=east,font=\footnotesize] at (-0.45,-1) {férié au Japon (mer.)};
\node[anchor=east,font=\footnotesize] at (-0.45,-2) {férié aux États-Unis (mer.)};
\node[anchor=east,font=\footnotesize] at (-0.45,-3) {férié aux États-Unis (jeu.)};
\foreach \r/\v/\h in {0/2/-1,-1/3/1,-2/2/1,-3/3/2} {
  \foreach \d in {0,1,2,3} {\draw[gray!50] (\d-0.4,\r-0.35) rectangle (\d+0.4,\r+0.35);}
  \node[omDef,font=\scriptsize] at (0,\r) {transaction};
  \fill[omThm!25] (\v-0.4,\r-0.35) rectangle (\v+0.4,\r+0.35);
  \node[omThm,font=\footnotesize\bfseries] at (\v,\r) {valeur};
}
\node[omProp,font=\scriptsize] at (1,-1) {JP fermé};
\node[omProp,font=\scriptsize] at (1,-2) {US fermés};
\node[omProp,font=\scriptsize] at (2,-3) {US fermés};
\end{tikzpicture}

\omcaption{\omterm{def:m2:the-fx-market:spot}{Dates de valeur spot} d'une transaction dollar--yen conclue un mardi. Un jour férié japonais
intermédiaire repousse le règlement au vendredi~; un jour férié américain ce même jour
ne le fait pas~; un jour férié américain à la date de valeur, en revanche, le fait. La règle telle que décrite par
les économistes de la Réserve fédérale en 2003~; jours fériés
fictifs.}\label{fig:m2:the-fx-market:spotdate}
\end{omfigure}

Un \omterm{def:m2:the-fx-market:spot}{cours croisé} peut toujours être construit à partir des deux paires en dollars, et la
construction fixe le côté retenu pour chaque jambe.

\begin{proposition}[Cours croisé par triangulation]\label{prop:m2:the-fx-market:cross}
Avec EURUSD coté $b_1/a_1$ et USDJPY $b_2/a_2$, le cours synthétique EURJPY vaut
$b_1b_2 / a_1a_2$. Avec EURUSD $b_1/a_1$ et GBPUSD $b_3/a_3$, le cours synthétique
EURGBP vaut $(b_1/a_3)/(a_1/b_3)$. Dans chaque cas, la fourchette relative du \omterm{def:m2:the-fx-market:spot}{cours
croisé} est voisine de la somme des fourchettes relatives des deux jambes.
\end{proposition}

\begin{proof}
Pour vendre des euros contre des yens par l'intermédiaire du dollar, on vend les euros au cours acheteur EURUSD $b_1$ et
on vend ces dollars contre des yens au cours acheteur USDJPY $b_2$~: $b_1b_2$ yens par euro. Pour acheter
des euros, on paie les deux cours vendeurs. Pour EURGBP, vendre des euros rapporte $b_1$ dollars, et
acheter des livres avec des dollars coûte $a_3$ dollars par livre. Les fourchettes relatives s'additionnent
au premier ordre, car $\log(a_1a_2) - \log(b_1b_2) = (\log a_1 - \log b_1) +
(\log a_2 - \log b_2)$.
\end{proof}

\begin{example}[Trois cours croisés]\label{ex:m2:the-fx-market:crosses}
Avec EURUSD 1.1462/1.1464, USDJPY 156.86/156.88 et GBPUSD 1.3371/1.3373, le
cours synthétique EURJPY est de 179.7929/179.8472, avec une largeur de 5.43 \omterm{def:m2:the-fx-market:pair}{pips}, soit une fourchette relative de
3.02 points de base, contre 1.74 et 1.27 pour les jambes~; EURGBP est
0.85710/0.85738, 2.78 \omterm{def:m2:the-fx-market:pair}{pips}~; GBPJPY 209.7375/209.7956, 5.81 \omterm{def:m2:the-fx-market:pair}{pips}. Un dealer qui
cote directement un \omterm{def:m2:the-fx-market:spot}{cours croisé} peut le coter plus serré que le cours synthétique, mais pas
croisé avec lui~: un cours acheteur direct sur EURJPY au-dessus de 179.8472 pourrait être vendu
puis racheté à travers les jambes, avec un profit.
\end{example}

\section{Dealers, niveaux et structure de gré à gré}

Le marché a été bâti par des banques qui négociaient entre elles et avec leurs clients,
et il en garde encore la forme. L'enquête 2025 compte 46~\% du volume négocié comme transactions
entre dealers déclarants et 50~\% comme transactions avec d'autres institutions financières,
telles que des banques plus petites, des sociétés de gestion, des hedge funds et des sociétés de trading agissant en principal~;
les entreprises non financières représentent le reste. Quatre places, le Royaume-Uni, les
États-Unis, Singapour et Hong Kong, concentrent les trois quarts du volume, Londres à elle seule
38~\%.

Un grand dealer cote des milliers de clients sur des centaines de paires, et compense leurs
flux les uns contre les autres avant de les porter sur le marché~: un client qui vend des euros
et un autre qui les achète dans la même seconde ne laissent au dealer aucun risque,
et aucune de ces deux transactions n'atteint une plateforme entre dealers. Cette internalisation (livre~1, chapitre~10) explique pourquoi le marché des dealers est concentré~:
plus un dealer voit de flux,
plus il compense, plus ses cotations sont serrées et plus il attire de flux.
Les banques plus petites et les dealers régionaux, qui forment le second rang, prennent les prix
des grands dealers et cotent leurs propres clients en s'y appuyant. Les sociétés de trading agissant en principal, qui négocient
avec leurs propres fonds au moyen d'algorithmes rapides, sont devenues elles aussi des teneurs de marché,
diffusant des prix aux banques et aux plateformes (\cref{ch:m2:fx-spot-microstructure}).

\begin{omfigure}
\centering
\begin{tikzpicture}[font=\small,
  box/.style={draw,thick,rounded corners=2pt,align=center,minimum height=0.9cm,minimum width=2.6cm,inner sep=3pt},
  arr/.style={{Stealth[length=2mm]}-{Stealth[length=2mm]},thick}]
\node[box,draw=omThm,fill=omThm!10] (pv) at (5,3.2) {plateformes principales\\(carnets d'ordres à cours limité centraux)};
\node[box,draw=omDef,fill=omDef!8] (d1) at (1.6,1.4) {grands dealers};
\node[box,draw=omDef,fill=omDef!8] (ptf) at (8.4,1.4) {sociétés de trading\\agissant en principal};
\node[box,draw=gray,fill=gray!8] (ecn) at (5,0.2) {ECN et\\plateformes multi-dealers};
\node[box,draw=omProp,fill=omProp!8] (sdp) at (1.6,-1.4) {plateformes\\mono-dealer};
\node[box,draw=omProp,fill=omProp!8] (cl) at (5,-3.0) {clients~: fonds, entreprises,\\banques plus petites};
\node[box,draw=omThm,fill=omThm!8] (pb) at (8.4,-1.4) {prime broker\\(crédit)};
\draw[arr] (d1) -- (pv); \draw[arr] (ptf) -- (pv);
\draw[arr] (d1) -- (ecn); \draw[arr] (ptf) -- (ecn);
\draw[arr] (d1) -- (sdp); \draw[arr] (sdp) -- (cl);
\draw[arr] (ecn) -- (cl);
\draw[arr,dashed] (cl) -- (pb); \draw[arr,dashed] (pb) -- (ptf);
\end{tikzpicture}

\omcaption{La structure du marché spot. Dealers et sociétés de trading agissant en principal
échangent entre eux sur les \omterm{def:m2:the-fx-market:venue}{plateformes principales} et sur d'autres réseaux électroniques, et avec
leurs clients sur ces réseaux et sur leurs propres \omterm{def:m2:the-fx-market:sdp}{plateformes mono-dealer}~;
un prime broker prête à un client son nom et son crédit pour négocier avec des
dealers qu'il ne pourrait pas affronter autrement (en pointillés). Schéma simplifié, d'après les descriptions de la Bank for
International Settlements.}\label{fig:m2:the-fx-market:structure}
\end{omfigure}

\section{Plateformes principales et réseaux électroniques}

\begin{definition}[Plateforme principale, réseau de communication électronique]\label{def:m2:the-fx-market:venue}
Une \emph{plateforme principale}\index{plateforme principale} est l'un des carnets d'ordres électroniques entre dealers
que le marché retient comme référence pour le prix d'une paire. Un
\emph{réseau de communication électronique}\index{réseau de communication électronique}
(ECN) est toute autre plateforme électronique de négociation où plusieurs fournisseurs de
liquidité et preneurs de liquidité se rencontrent, par carnet d'ordres ou par demande de cotation (RFQ), anonymement
ou avec les noms divulgués.
\end{definition}

Depuis trois décennies, deux courtiers électroniques, Refinitiv Matching, anciennement
Reuters, et EBS Market, aujourd'hui détenu par CME Group, sont les \omterm{def:m2:the-fx-market:venue}{plateformes principales} du
marché spot,
carnets d'ordres à cours limité centraux sur lesquels des dealers, et depuis le milieu des années 2000 des sociétés de trading agissant en principal,
affichent et frappent des prix~; leurs prix ont servi de référence pour
le reste du marché. Le BIS a compté plus de trente plateformes secondaires dans le
segment dealer-client en 2022, et les volumes sur les \omterm{def:m2:the-fx-market:venue}{plateformes principales} en recul
depuis plus d'une décennie~: les dealers internalisent davantage, les algorithmes d'exécution découpent
les ordres entre de nombreuses plateformes, et des dealers craignant d'être pris à revers par des sociétés plus rapides
ont poussé les plateformes à ajouter aux ordres entrants de petits délais, des speed bumps.
Une part importante de la formation des prix se fait désormais aussi sur les contrats à terme sur devises.

\section{Plateformes mono-dealer et prime brokerage}

\begin{definition}[Plateforme mono-dealer, prime brokerage de change]\label{def:m2:the-fx-market:sdp}
Une \emph{plateforme mono-dealer}\index{plateforme mono-dealer} est le système électronique propre à un dealer
par lequel ses clients voient ses prix et négocient avec lui.
Le \emph{prime brokerage de change}\index{prime brokerage de change} est le service par lequel un
dealer permet à un client de négocier avec un ensemble d'autres dealers au nom du prime broker
et sur son crédit~: chaque transaction fait l'objet d'un give-up auprès du prime broker, qui devient
la contrepartie du dealer exécutant et du client.
\end{definition}

L'enquête 2022 a enregistré un basculement de sept points de pourcentage des négociations vers
les canaux électroniques directs, les \omterm{def:m2:the-fx-market:sdp}{plateformes mono-dealer} et les flux de prix directs,
en s'éloignant des plateformes anonymes. Pour un client, une \omterm{def:m2:the-fx-market:sdp}{plateforme mono-dealer} offre
des prix personnalisés et, souvent, les autres services du dealer~; pour un dealer, c'est
le meilleur moyen de voir les flux de ses clients et de les internaliser.

Le prime brokerage permet à un hedge fund ou à une société de trading agissant en principal
d'agir comme une banque. Le dealer exécutant ne voit que le nom du prime broker, si bien que le
client négocie de façon anonyme et avec de nombreux dealers sous une seule relation de crédit, un seul
ensemble de règlements-livraisons et un seul rapport~; le prime broker perçoit une commission et supporte le
risque de crédit du client. Le BIS a estimé que près d'un tiers du volume
des transactions avec des clients financiers passait par le prime brokerage en 2019, et voir au \cref{ch:m2:getting-access-fx}
ce qu'il faut pour obtenir un tel arrangement, ainsi que
les limites qui s'y rattachent.

\section{Ce que disent les enquêtes sur les volumes}

Tous les trois ans, les banques centrales demandent aux dealers de leur ressort de
déclarer leurs opérations d'avril, et le BIS en rassemble les résultats~: en 2025, 52
juridictions et plus de 1\,100 dealers. L'enquête compte chaque transaction une seule fois
à l'échelle mondiale, selon le lieu du desk de vente du dealer.

\begin{omfigure}
\begin{tikzpicture}
\begin{axis}[width=11cm,height=4.8cm,ybar,bar width=10pt,
  ylabel={part du volume (\%)},
  tick label style={font=\small},label style={font=\small},
  xmin=-0.6,xmax=4.6,ymin=0,ymax=60,grid=major,grid style={gray!25},
  xtick={0,1,2,3,4},xticklabels={spot,change à\\terme sec,swaps de\\change,options,swaps de\\devises},
  xticklabel style={align=center},
  nodes near coords,every node near coord/.append style={font=\scriptsize},
  legend columns=2,legend style={font=\footnotesize,at={(0.5,-0.42)},anchor=north,draw=none,fill=none,/tikz/every even column/.append style={column sep=8pt}},
  table/col sep=comma]
\addplot[fill=gray!40,draw=gray] table[x=k,y=y2022]{figdata/markets-2/14-the-fx-market/instruments.csv};
\addplot[fill=omThm!60,draw=omThm] table[x=k,y=y2025]{figdata/markets-2/14-the-fx-market/instruments.csv};
\legend{avril 2022 (7.5~billions de dollars par jour),avril 2025 (9.6~billions de dollars par jour)}
\end{axis}
\end{tikzpicture}

\omcaption{Part du volume mondial de change par instrument, en pourcentage du total quotidien, en
avril 2022 et avril 2025. Les swaps de change restent l'instrument le plus important, mais le spot,
les changes à terme sec et les options ont progressé plus vite. Données~: enquête triennale de la BRI auprès des banques centrales,
édition 2025.}\label{fig:m2:the-fx-market:instruments}
\end{omfigure}

Deux constats se dégagent. D'abord, le spot pèse moins d'un tiers du marché~: l'essentiel
de la négociation porte sur les swaps de change et les changes à terme, instruments de financement et de couverture
(\cref{fig:m2:the-fx-market:instruments}). Ensuite, le dollar est le
véhicule du marché (\cref{fig:m2:the-fx-market:currencies})~: les dix paires les plus
échangées le comprennent toutes, et un exportateur suédois payé en yens trouve souvent
moins coûteux de vendre ses yens contre des dollars, puis ses dollars contre des couronnes, que de traiter le
\omterm{def:m2:the-fx-market:spot}{cours croisé}, peu liquide.
Les parts totalisent 200~\% car chaque transaction implique deux devises. L'enquête
signale aussi qu'avril 2025 n'a pas été un mois calme~: il a suivi des annonces de politique
commerciale qui ont fait fortement bouger les devises.

\begin{omfigure}
\begin{tikzpicture}
\begin{axis}[width=11cm,height=4.6cm,ybar,bar width=14pt,
  ylabel={part des transactions (\%)},
  tick label style={font=\small},label style={font=\small},
  xmin=-0.6,xmax=6.6,ymin=0,ymax=110,grid=major,grid style={gray!25},ytick={0,20,40,60,80,100},
  xtick={0,1,2,3,4,5,6},xticklabels={USD,EUR,JPY,GBP,CNY,CHF,HKD},
  nodes near coords,every node near coord/.append style={font=\scriptsize},
  table/col sep=comma]
\addplot[fill=omDef!60,draw=omDef] table[x=k,y=y2025]{figdata/markets-2/14-the-fx-market/currencies.csv};
\end{axis}
\end{tikzpicture}

\omcaption{Part des transactions de change mondiales impliquant chaque devise, avril 2025 (les
parts totalisent 200~\% puisque chaque transaction a deux côtés). Les dollars australien et
canadien, à environ 6~\% chacun, ne sont pas représentés. Données~: enquête triennale de la BRI
auprès des banques centrales, édition 2025.}\label{fig:m2:the-fx-market:currencies}
\end{omfigure}

\begin{dated}{2026-09}{Taille et règlement-livraison}\label{dat:m2:the-fx-market:survey}
Enquête triennale de la BRI, avril 2025~: 9.6~billions de dollars par jour~; spot 3~billions de dollars,
changes à terme sec 1.8~billion de dollars, swaps de change 4~billions de dollars~; le dollar figure d'un côté
de 89.2~\% des transactions, l'euro 28.9~\%, le yen 16.8~\%, la livre sterling 10.2~\%, le renminbi 8.5~\%,
le franc suisse 6.4~\%. L'enquête est répétée tous les trois ans. CLS, le système de règlement-livraison
des principales devises (\cref{ch:m2:settlement-risk}), règle 18
devises, en moyenne plus de 8~billions de dollars par jour, pour plus de 75
membres du règlement-livraison.
\end{dated}

\section{Tutoriel~: dates spot et cours croisés}\label{tut:m2:the-fx-market:pairs}

\begin{tutorial}
\textbf{Objectif.} Nommer les paires selon la convention, calculer les \omterm{def:m2:the-fx-market:spot}{dates de valeur spot} à travers deux
calendriers de jours fériés, et trianguler les \omterm{def:m2:the-fx-market:spot}{cours croisés} à partir de leurs cours acheteur et vendeur.
\textbf{État final~:} \cref{fig:m2:the-fx-market:spotdate},
\cref{ex:m2:the-fx-market:crosses} et les chiffres du problème du week-end.
\begin{tutsteps}
\item \textbf{La date spot}, avec l'exception du dollar le jour
intermédiaire.
\omcode{code/firm/fxpairs/firm_fxpairs.py}{35}{50}{Date de valeur spot d'une paire à travers deux calendriers de jours fériés.}\label{lst:m2:the-fx-market:spot}
\item \textbf{Le \omterm{def:m2:the-fx-market:spot}{cours croisé}}, chaque jambe du côté imposé par la transaction, et le
contrôle d'arbitrage d'une cotation directe.
\omcode{code/firm/fxpairs/firm_fxpairs.py}{53}{81}{Cours croisé via le dollar, et test d'arbitrage.}\label{lst:m2:the-fx-market:cross}
\item \textbf{Exécuter} \texttt{fx\_demo.crosses()}, \texttt{fx\_demo.spot\_examples()},
\texttt{fx\_demo.the\_cross()} et \texttt{fig\_fx.py}.
\end{tutsteps}
\textbf{Pour aller plus loin.} Ajouter un calendrier de jours fériés issu d'une source publique pour
la zone euro (les jours de fermeture de TARGET) et le Japon, puis calculer une année
de dates spot EURJPY~; lister ensuite les jours où la date spot saute de plus de deux
jours ouvrés.
\end{tutorial}

\section{Construction~: les conventions de paires}\label{bld:m2:the-fx-market:fxpairs}

\begin{build}
\bfield{Rôle} Tout système de change de la mini-société commence ici~: le nom
et l'orientation d'une paire, son \omterm{def:m2:the-fx-market:pair}{pip}, sa date spot et le prix synthétique
d'un \omterm{def:m2:the-fx-market:spot}{cours croisé}.
\bfield{Interface} \texttt{pair\_name(a, b, priority)}~; \texttt{pip(pair)}~;
\texttt{spot\_date(trade, pair, holidays)}~; \texttt{Quote(bid, ask)}~;
\texttt{cross\_via\_usd(leg1, leg2, target)}~; \texttt{spread\_pips}~;
\texttt{arbitrage(direct, synthetic)}.
\bfield{Règles} Priorité de la \omterm{def:m2:the-fx-market:pair}{devise de base} conservée comme donnée~; jours fériés passés en paramètre, jamais
codés en dur~; T+2 sur les jours ouvrés des deux devises, T+1 pour USDCAD, un jour férié américain le
jour intermédiaire ignoré pour les paires en dollar~; chaque jambe d'un \omterm{def:m2:the-fx-market:spot}{cours croisé} du côté imposé par la
transaction.
\bfield{Tests d'acceptation} \texttt{code/firm/fxpairs/tests/}~: noms et \omterm{def:m2:the-fx-market:pair}{pips} des paires~;
les quatre cas publiés dollar--yen et USDCAD~; une transaction du vendredi réglée le
mardi~; des \omterm{def:m2:the-fx-market:spot}{cours croisés} égaux aux produits et quotients des bons côtés~; une
cotation directe croisée signalée comme arbitrage.
\bfield{Pour aller plus loin} Calendriers de jours fériés réels, versionnés~; le roll de 17~h à New York~;
les dates de valeur des changes à terme (\cref{ch:m2:fx-swaps-forwards-and-the-cross-currency-basis})~; les \omterm{def:m2:the-fx-market:spot}{cours
croisés} passant par une devise autre que le dollar.
\end{build}

\begin{omsources}
\item Bank for International Settlements, Triennial Central Bank Survey,
«~OTC foreign exchange turnover in April 2025~» et communiqué de presse, 30 septembre
2025.
\item Bank for International Settlements, \emph{Quarterly Review}, décembre 2019
(«~Sizing up global foreign exchange markets~») et décembre 2022 (encadrés sur
l'exécution des transactions et le marché spot fragmenté).
\item A.~Chaboud et J.~Wright, «~Uncovered interest parity: it works, but not
for long~», Federal Reserve Board, International Finance Discussion Paper 752,
2003.
\item Federal Reserve, H.10 foreign exchange rates~; instructions de l'enquête FR 3036~;
Foreign Exchange Committee, amendements aux 1998 FX
definitions, 2011.
\item CLS Group, informations sur la société.
\end{omsources}

\section{Exercices}

\begin{exercise}[$\star$]\label{exo:m2:the-fx-market:1}
EURUSD vaut 1.1462/1.1464. Un client vend 5~millions d'euros. Combien de dollars
reçoit-il, et quelle est la fourchette en \omterm{def:m2:the-fx-market:pair}{pips} et en points de base du cours médian~?
\end{exercise}

\begin{exercise}[$\star$]\label{exo:m2:the-fx-market:2}
Nommer les paires formées par le dollar et le yen, l'euro et la livre sterling, le
dollar australien et le yen, et donner le \omterm{def:m2:the-fx-market:pair}{pip} de chacune.
\end{exercise}

\begin{exercise}[$\star$]\label{exo:m2:the-fx-market:3}
Une transaction USDJPY est conclue le mardi 6 octobre 2026. Donner sa \omterm{def:m2:the-fx-market:spot}{date de valeur spot} sans
jour férié, avec un jour férié japonais le mercredi, et avec un jour férié américain le
jeudi. Et pour une transaction USDCAD le même jour~?
\end{exercise}

\begin{exercise}[$\star\star$]\label{exo:m2:the-fx-market:4}
À partir de EURUSD 1.1462/1.1464 et de GBPUSD 1.3371/1.3373, déduire le cours synthétique
de l'EURGBP et sa fourchette en \omterm{def:m2:the-fx-market:pair}{pips}.
\end{exercise}

\begin{exercise}[$\star\star$]\label{exo:m2:the-fx-market:5}
Pourquoi les parts de devises de l'enquête totalisent-elles 200~\%, et que dit la part de
89~\% du dollar sur la manière dont un \omterm{def:m2:the-fx-market:spot}{cours croisé} comme SEKJPY est habituellement négocié~?
\end{exercise}

\begin{exercise}[$\star\star$]\label{exo:m2:the-fx-market:6}
Un hedge fund ayant un prime broker exécute une transaction de 50~millions de dollars sur
EURUSD auprès d'un dealer avec lequel il n'a aucun accord. Qui est la contrepartie du dealer,
et qui supporte le risque de crédit du fonds~?
\end{exercise}

\begin{exercise}[$\star\star\star$]\label{exo:m2:the-fx-market:7}
\emph{Programmation.} Avec \texttt{cross\_via\_usd}, donner le GBPJPY synthétique et sa
fourchette relative, et vérifier qu'elle est proche de la somme des fourchettes relatives des
jambes.
\end{exercise}

\begin{exercise}[$\star\star\star$]\label{exo:m2:the-fx-market:8}
\emph{Repérer l'erreur.} «~Le volume quotidien du marché des changes est de 9.6~billions de dollars, donc un exportateur peut
vendre n'importe quel montant de n'importe quelle devise au cours affiché.~» Corriger-la.
\end{exercise}

\section{Problème~: le cours croisé}

\begin{problem}[{Problème du week-end --- le montant maximal qu'un client doit payer pour un cours croisé}]\label{pb:m2:the-fx-market:1}
Un client veut acheter 10~millions d'euros contre des yens. Le marché affiche EURUSD
1.1462/1.1464 et USDJPY 156.86/156.88. Deux dealers cotent EURJPY directement~:
dealer~A 179.80/179.84, dealer~B 179.85/179.87.

\medskip\noindent\textbf{Partie I --- Le \omterm{def:m2:the-fx-market:spot}{cours croisé} synthétique.}
\begin{enumerate}
\item Quel cours de chaque jambe le client consomme-t-il pour acheter des euros contre des yens~?
\item Donner le cours acheteur et le cours vendeur synthétiques de EURJPY.
\item Donner sa fourchette en \omterm{def:m2:the-fx-market:pair}{pips} et en points de base du cours médian.
\item Donner les fourchettes relatives des jambes et comparer leur somme à celle du \omterm{def:m2:the-fx-market:spot}{cours croisé}.
\item Combien de dollars transitent par la transaction si le client achète de manière synthétique~?
\end{enumerate}

\medskip\noindent\textbf{Partie II --- Les cotations directes.}
\begin{enumerate}[resume]
\item La cotation du dealer~A est-elle cohérente avec les jambes~?
\item Et celle du dealer~B~? Que ferait un arbitragiste~?
\item Donner le profit d'arbitrage sur 10~millions d'euros, en yens et en dollars.
\item Pourquoi une telle cotation peut-elle exister un instant, et pourquoi pas longtemps~?
\item Où le client doit-il acheter~?
\end{enumerate}

\medskip\noindent\textbf{Partie III --- Coûts au-delà de la cotation.}
\begin{enumerate}[resume]
\item Pourquoi négocier les jambes coûte-t-il plus cher au client que ne le laisse supposer le calcul~?
\item Quand le \omterm{def:m2:the-fx-market:spot}{cours croisé} direct est-il moins cher que le synthétique~?
\item En quoi le règlement-livraison diffère-t-il entre les deux circuits~?
\item Qu'est-ce qui change si le mercredi est férié au Japon et si la transaction a lieu le mardi~?
\item Pourquoi un dealer peut-il coter le \omterm{def:m2:the-fx-market:spot}{cours croisé} avec une fourchette plus serrée que le synthétique~?
\end{enumerate}

\medskip\noindent\textbf{Partie IV --- Jugement.}
\begin{enumerate}[resume]
\item Quel est le danger de coter le \omterm{def:m2:the-fx-market:spot}{cours croisé} à partir de jambes périmées~?
\item Pourquoi les \omterm{def:m2:the-fx-market:spot}{cours croisés} peu liquides passent-ils par le dollar~?
\item Comment un ordre important sur le \omterm{def:m2:the-fx-market:spot}{cours croisé} ferait-il bouger les jambes~?
\item Énoncer le \emph{résultat nommé}~: la fourchette la plus large que le client doit accepter
pour EURJPY aujourd'hui.
\item En une phrase~: qu'est-ce qui borne le prix de n'importe quel \omterm{def:m2:the-fx-market:spot}{cours croisé}~?
\end{enumerate}
\end{problem}

\section{Questions d'entretien}

\begin{interviewq}[$\star$ \iqroles{trader, développeur}]\label{iq:m2:the-fx-market:1}
Que signifie EURUSD 1.1462/1.1464, et quel cours le client qui achète des euros
consomme-t-il~?
\end{interviewq}

\begin{interviewq}[$\star$ \iqroles{développeur}]\label{iq:m2:the-fx-market:2}
Comment calculeriez-vous la \omterm{def:m2:the-fx-market:spot}{date de valeur spot} d'une transaction de change~? Quels cas limites testeriez-vous~?
\end{interviewq}

\begin{interviewq}[$\star\star$ \iqroles{trader, chercheur}]\label{iq:m2:the-fx-market:3}
Déduire le cours acheteur et le cours vendeur de EURGBP à partir de EURUSD et de GBPUSD.
\end{interviewq}

\begin{interviewq}[$\star\star$ \iqroles{trader, banque}]\label{iq:m2:the-fx-market:4}
Pourquoi le marché des changes est-il de gré à gré plutôt que coté en bourse, et comment le
trading électronique l'a-t-il transformé~?
\end{interviewq}

\begin{interviewq}[$\star\star$ \iqroles{banque}]\label{iq:m2:the-fx-market:5}
Que fait un prime broker pour un client de change, et quels risques prend-il~?
\end{interviewq}

\begin{interviewq}[$\star\star\star$ \iqroles{développeur, trader}]\label{iq:m2:the-fx-market:6}
Concevoir le composant qui publie les prix synthétiques d'une centaine de
\omterm{def:m2:the-fx-market:spot}{cours croisés} à partir des jambes en dollars, avec un budget de latence de quelques microsecondes.
\end{interviewq}
